\documentclass[11pt]{article}
\usepackage[a4paper,margin=1in]{geometry}
\usepackage{amsmath,amssymb,amsthm,mathtools}
\usepackage{enumitem}
\usepackage{hyperref}

\title{\bf Various Channel Models in the Current Codebase}
\author{}
\date{}

\newtheorem{remark}{Remark}

\begin{document}
\maketitle

\begin{abstract}
This document gives a methods-level description of all physical noise channels
currently implemented in
\texttt{code/quantum\_simulation/noise\_engine.py} and
\texttt{code/quantum\_simulation/noise\_channels.py}.
For each channel, we provide: (i) mathematical construction,
(ii) parameter meaning and selection guidance, and
(iii) exact simulation semantics used by the code.
The target usage is direct reuse in paper Methods/Setup sections.
\end{abstract}

\tableofcontents

\section{Unified Simulation Setup}

We model the ideal circuit as an ordered instruction sequence.
Let instruction index be $i=1,\dots,L$, and let
$\rho_i$ be the state before instruction $i$.
The simulator applies
\begin{equation}
\rho_{i+1} = \mathcal{N}_i\!\left(\mathcal{U}_i(\rho_i)\right),
\qquad
\mathcal{U}_i(\rho)=U_i\rho U_i^\dagger,
\label{eq:generic_update}
\end{equation}
where $\mathcal{N}_i$ is channel-specific and can depend on control mismatch,
gate type, and idle-window position.

\paragraph{Timeline convention.}
The code builds a serial timeline with fixed idle gap $\tau_{\mathrm{idle}}$
between adjacent instructions.
Thus each circuit has $(L-1)$ idle windows:
\begin{equation}
W_i = [t_i^{\mathrm{end}},\, t_i^{\mathrm{end}}+\tau_{\mathrm{idle}}],
\qquad i=1,\dots,L-1.
\end{equation}
Idle-only channels act on these windows, not on the final post-instruction gap.

\paragraph{Control mismatch notation.}
Let current control vector be $u\in\mathbb{R}^{d}$ and drifted optimum be
$u^\star(t)\in\mathbb{R}^{d}$.
Define global mismatch
\begin{equation}
m(t)=\frac{1}{d}\sum_{k=1}^{d}\left(u_k-u_k^\star(t)\right)^2.
\label{eq:global_mismatch}
\end{equation}

\paragraph{Probability clipping.}
Most channels use
\begin{equation}
\mathrm{clip}_{p_{\max}}(x)=\min\{p_{\max},\max\{0,x\}\},
\end{equation}
with $p_{\max}=\texttt{p\_clip\_max}$.

\section{Channel I: Google-Like Global Gate Depolarizing (\texttt{google\_global})}

\subsection{Channel construction}
This channel adds depolarizing noise after each non-measure/non-reset gate.
The same mismatch scalar $m(t)$ drives all gates at time $t$:
\begin{align}
p_{1q}(t) &= \mathrm{clip}_{p_{\max}}\!\left(p_{1q}^{0}+s_{1q}\,m(t)\right),\\
p_{2q}(t) &= \mathrm{clip}_{p_{\max}}\!\left(p_{2q}^{0}+s_{2q}\,m(t)\right).
\label{eq:global_google_probs}
\end{align}
Here $p_{1q}^{0},p_{2q}^{0}$ are base rates and $s_{1q},s_{2q}$ are
sensitivity coefficients.

For a 1-qubit gate, code appends \texttt{DEPOLARIZE1($p_{1q}$)};
for a 2-qubit gate, \texttt{DEPOLARIZE2($p_{2q}$)}.
Stim semantics are uniform Pauli mixtures:
\begin{align}
\mathcal{D}_{1}(p)(\rho)&=(1-p)\rho+\frac{p}{3}\sum_{P\in\{X,Y,Z\}}P\rho P,\\
\mathcal{D}_{2}(p)(\rho)&=(1-p)\rho+\frac{p}{15}\sum_{P\in\mathcal{P}_2\setminus\{II\}}P\rho P.
\end{align}

\subsection{Parameter selection}
A practical calibration path is:
\begin{enumerate}[label=\textbf{G\arabic*}.]
\item Fix $p_{1q}^{0},p_{2q}^{0}$ from hardware baseline characterization.
\item Estimate expected mismatch scale
$\sigma_m\approx \mathbb{E}[m(t)]$ in the intended RL/control regime.
\item Choose sensitivities from desired dynamic range:
\begin{equation}
s_{1q}\approx \frac{\Delta p_{1q}}{\sigma_m},
\qquad
s_{2q}\approx \frac{\Delta p_{2q}}{\sigma_m},
\end{equation}
where $\Delta p$ is the target increase from baseline at typical mismatch.
\item Set $p_{\max}$ large enough not to clip normal operation,
but low enough to avoid unphysical saturation.
\end{enumerate}

\section{Channel II: Google-Like Gate-Specific Depolarizing (\texttt{google\_gate\_specific})}

\subsection{Channel construction}
This model preserves the same depolarizing insertion rule as Channel I,
but uses gate-specific mismatch features.
Control is partitioned as
\begin{equation}
u=(u^{(1q)},u^{(2q)}),
\quad
u^{(1q)}\in\mathbb{R}^{N_{1q}},
\quad
u^{(2q)}\in\mathbb{R}^{N_{2q}}.
\end{equation}
Each instruction is deterministically assigned a slot index
(using stable hash of gate name + qubit targets):
\begin{equation}
j_{1q}=h_{1q}(\text{inst}),
\qquad
j_{2q}=h_{2q}(\text{inst}).
\end{equation}
Per-instruction probabilities are
\begin{align}
p_{1q,i}&=\mathrm{clip}_{p_{\max}}\!\left(
p_{1q}^{0}+s_{1q}\left(u^{(1q)}_{j_{1q}}-u^{\star(1q)}_{j_{1q}}(t_i)\right)^2
\right),\\
p_{2q,i}&=\mathrm{clip}_{p_{\max}}\!\left(
p_{2q}^{0}+s_{2q}\left(u^{(2q)}_{j_{2q}}-u^{\star(2q)}_{j_{2q}}(t_i)\right)^2
\right).
\label{eq:gate_specific_probs}
\end{align}

\subsection{Parameter selection}
\begin{enumerate}[label=\textbf{S\arabic*}.]
\item Choose $N_{1q},N_{2q}$ to control parameter sharing.
Larger slot counts reduce hash collisions; smaller counts enforce sharing.
\item Keep $p_{1q}^{0},p_{2q}^{0}$ at calibrated baselines,
then tune $s_{1q},s_{2q}$ for desired control leverage.
\item Verify collision statistics for the target circuit family,
since collision structure changes effective model capacity.
\end{enumerate}

\begin{remark}
With one slot ($N_{1q}=N_{2q}=1$), this channel degenerates toward global behavior.
With many slots, it approximates gate-by-gate control sensitivity.
\end{remark}

\section{Channel III: Idle Depolarizing Pauli Channel (\texttt{idle\_depolarizing})}

\subsection{Channel construction}
This is an action-independent, idle-window-only Pauli channel.
Given per-idle total probability $p_{\mathrm{idle}}$ and axis weights
$w_x,w_y,w_z\ge 0$, define normalized weights
\begin{equation}
\tilde w_\alpha = \frac{w_\alpha}{w_x+w_y+w_z},
\qquad \alpha\in\{x,y,z\}.
\end{equation}
Then per-idle axis probabilities are
\begin{equation}
p_x=\tilde w_x p_{\mathrm{idle}},
\quad
p_y=\tilde w_y p_{\mathrm{idle}},
\quad
p_z=\tilde w_z p_{\mathrm{idle}}.
\label{eq:idle_axis_probs}
\end{equation}
In code, this is implemented via constant rate functions
$\lambda_\alpha=p_\alpha/\tau_{\mathrm{idle}}$ integrated over each idle window,
which exactly recovers Eq.~\eqref{eq:idle_axis_probs}.

At each idle window and each qubit, simulator appends
\texttt{PAULI\_CHANNEL\_1([$p_x,p_y,p_z$])}.

\subsection{Parameter selection}
\begin{enumerate}[label=\textbf{I\arabic*}.]
\item Choose $p_{\mathrm{idle}}$ from target per-window error budget.
\item Use equal weights for isotropic idle noise; use biased weights
for dephasing-dominant or amplitude-like regimes.
\item Ensure validity: $p_x,p_y,p_z\ge 0$ and
$p_x+p_y+p_z\le 1$.
\item If starting from a per-circuit budget $P_{\mathrm{circ}}$ over $N_w$ idle windows,
convert by
\begin{equation}
p_{\mathrm{idle}}=1-(1-P_{\mathrm{circ}})^{1/N_w}.
\label{eq:idle_budget_conversion}
\end{equation}
\end{enumerate}

\section{Channel IV: Parametric Google Gate-Specific (\texttt{parametric\_google})}

\subsection{Channel construction}
This channel uses the same gate-specific machinery as Channel II,
but introduces explicit regime parameters
$a=\texttt{channel\_regime\_a}$ and
$b=\texttt{channel\_regime\_b}$.
Code rescales both base rates and sensitivities:
\begin{align}
p_{1q}^{0}&\mapsto a\,p_{1q}^{0}, & s_{1q}&\mapsto a\,s_{1q},\\
p_{2q}^{0}&\mapsto b\,p_{2q}^{0}, & s_{2q}&\mapsto b\,s_{2q}.
\end{align}
Hence Eq.~\eqref{eq:gate_specific_probs} becomes
\begin{align}
p_{1q,i}&=\mathrm{clip}_{p_{\max}}\!\left(a\left[p_{1q}^{0}+s_{1q}\delta_{1q,i}^2\right]\right),\\
p_{2q,i}&=\mathrm{clip}_{p_{\max}}\!\left(b\left[p_{2q}^{0}+s_{2q}\delta_{2q,i}^2\right]\right),
\end{align}
where $\delta_{1q,i},\delta_{2q,i}$ denote corresponding slot mismatch magnitudes.

\subsection{Parameter selection}
\begin{enumerate}[label=\textbf{P\arabic*}.]
\item Use $(a,b)$ for structured regime sweeps while keeping architecture fixed.
\item $a$ primarily controls 1q channel severity; $b$ controls 2q severity.
\item For fair comparisons, keep all non-channel hyperparameters fixed and vary only $(a,b)$.
\end{enumerate}

\section{Channel V: Correlated Pauli Idle Channel (\texttt{correlated\_pauli\_noise\_channel})}

\subsection{Design goal and constraints}
This is the only stateful, temporally correlated channel in current code.
It is built as a Hidden-Markov Pauli process injected on idle windows.
The implementation enforces:
\begin{enumerate}[label=\textbf{C\arabic*}.]
\item Qubit independence: each qubit has its own hidden chain(s).
\item Direction independence: X/Y/Z use identical $(f,g)$ statistics.
\item Equal axis weights are required ($w_x=w_y=w_z$).
\end{enumerate}

\subsection{Step 1: mismatch to target total probability}
Using global mismatch at current step,
\begin{equation}
m=\frac{1}{d}\sum_{k=1}^{d}(u_k-u_k^\star(0))^2,
\end{equation}
define base strength
\begin{equation}
s_{\mathrm{base}}=p_{1q}^{0}+s_{1q}m,
\end{equation}
and raw target
\begin{equation}
p_{\mathrm{target}}=\mathrm{clip}_{p_{\max}}(g\,s_{\mathrm{base}}).
\label{eq:corr_target}
\end{equation}

Two $g$ modes:
\begin{itemize}
\item \texttt{per\_window}: $p_{\mathrm{nom}}=p_{\mathrm{target}}$.
\item \texttt{per\_circuit}: for normalization windows $N_w$,
\begin{equation}
p_{\mathrm{nom}}=1-(1-p_{\mathrm{target}})^{1/N_w}.
\label{eq:per_circuit_norm}
\end{equation}
\end{itemize}
Then clip $p_{\mathrm{nom}}$ to $[0,p_{\max}]$.

\subsection{Step 2: convert total probability to per-axis mean}
Let $p$ be per-axis Bernoulli mean (same for X/Y/Z).
The non-identity probability in one idle window is
\begin{equation}
p_{\mathrm{tot}} = 1-\Pr(I)=1-\bigl((1-p)^3+p^3\bigr)=3p(1-p).
\label{eq:tot_to_axis}
\end{equation}
Code solves this as
\begin{equation}
p=\frac{1-\sqrt{1-\frac{4}{3}p_{\mathrm{tot}}}}{2},
\qquad
p_{\mathrm{tot}}\le 0.75.
\end{equation}
Hence it clips
$p_{\mathrm{tot}}=\min\{p_{\mathrm{nom}},0.75-\varepsilon\}$.

\subsection{Step 3: Hidden-Markov telegraph process}
For each qubit $q$ and each axis $\alpha\in\{X,Y,Z\}$,
introduce hidden state $s_{\alpha,q,t}\in\{0,1\}$ with emission probs
\begin{equation}
p_{\mathrm{low}}=0,
\qquad
p_{\mathrm{high}}=\min\{2p,1\}.
\end{equation}
Transition matrix is symmetric:
\begin{equation}
T=
\begin{pmatrix}
1-\gamma & \gamma\\
\gamma & 1-\gamma
\end{pmatrix},
\qquad
\gamma=\frac{1-e^{-f\tau_{\mathrm{idle}}}}{2},
\label{eq:corr_transition}
\end{equation}
where $f=\texttt{corr\_frequency\_hz}$.
Initial distribution is $(0.5,0.5)$.

\subsection{Step 4: per-window sampling and injection}
For each idle window $t$:
\begin{enumerate}[label=\textbf{H\arabic*}.]
\item Sample events
$e_{\alpha,q,t}\sim \mathrm{Bernoulli}(p_{s_{\alpha,q,t}})$.
\item Advance hidden states with Eq.~\eqref{eq:corr_transition}.
\item Append Pauli gates to noisy circuit:
\texttt{X} on all qubits with $e_{X,q,t}=1$, similarly for \texttt{Y}, \texttt{Z}.
\end{enumerate}
Multiple axes may fire on one qubit in the same window.

\subsection{Physical interpretation and parameter selection}
\paragraph{Correlation frequency $f$.}
From Eq.~\eqref{eq:corr_transition}, one-step correlation factor is
$\rho=e^{-f\tau_{\mathrm{idle}}}$.
Thus physical correlation time is approximately
\begin{equation}
\tau_c\approx \frac{1}{f}.
\end{equation}
Small $f$ gives longer memory (low-frequency drift), large $f$ gives faster mixing.

\paragraph{Strength $g$.}
From Eq.~\eqref{eq:corr_target}, $g$ scales the mismatch-driven base strength.
For a target budget $P$ and estimated mismatch $m$, a first-order choice is
\begin{equation}
g \approx \frac{P}{p_{1q}^{0}+s_{1q}m}
\quad
\text{(before clipping effects).}
\end{equation}

\paragraph{$g$-mode choice.}
\begin{itemize}
\item Use \texttt{per\_window} when all compared experiments have similar circuit length.
\item Use \texttt{per\_circuit} when comparing different round/circuit lengths and you want length-normalized budgets.
\end{itemize}

\subsection{Statefulness and runtime semantics}
This model is stateful across consecutive \texttt{apply(...)} calls within one shot.
Therefore:
\begin{enumerate}[label=\textbf{R\arabic*}.]
\item apply-result cache is disabled by design;
\item hidden states are reset at shot start via \texttt{start\_shot()};
\item Steane adapter forces effective \texttt{shot\_workers=1} for this channel.
\end{enumerate}
These constraints preserve correct temporal-memory semantics.

\section{Recommended Methods Reporting Template}

For reproducibility, report at least:
\begin{enumerate}[leftmargin=2em]
\item Exact channel key and resolved mode (avoid only writing ``auto").
\item Base rates/sensitivities: $p_{1q}^{0},p_{2q}^{0},s_{1q},s_{2q},p_{\max}$.
\item For gate-specific models: $N_{1q},N_{2q}$ and slot-mapping policy.
\item For idle models: $\tau_{\mathrm{idle}}$, axis weights, and per-window or per-circuit budget definition.
\item For correlated model: $(f,g)$, $g$-mode, normalization windows $N_w$, and any clipping regime.
\item Runtime semantics affecting statistics: trace mode, shot workers, state reset policy.
\end{enumerate}

\end{document}
